This research examined how key reinforcement learning parameters, model architectures, and algorithmic choices influence the stability and effectiveness of LLM fine-tuning for query generation under the DeepRetrieval framework~\cite{jiang2025deepretrieval}. Across all experiments, the learning rate emerged as the single most critical factor: high values induced divergence and oscillatory KL behavior, while smaller rates stabilized optimization but slowed learning progress. Reward scaling and secondary hyperparameters such as the KL penalty or PPO epochs affected smoothness and exploration but did not resolve underlying instability. These results emphasize that careful learning-rate control is essential for stable RL fine-tuning.

Model architecture and scale also proved decisive. The LLaMA-3.2-3B model consistently achieved smoother training, higher baseline rewards, and lower loss fluctuations than the smaller Qwen2-0.5B model, underscoring the stabilizing effect of model capacity and richer pretraining. However, stability gains did not translate into proportional learning improvements, revealing that optimization remained limited by the reward signal rather than model expressiveness.

Switching from PPO to GRPO further improved optimization dynamics, producing smoother KL divergence and faster, more stable learning of structured Boolean-style query formats. While GRPO enhanced efficiency, recall-based learning remained challenging, suggesting that the sparse and coarse nature of the recall@50 reward limited the model’s ability to learn retrieval-oriented behaviors. Using higher-coverage metrics such as recall@3000 could strengthen the reward signal and accelerate learning in future studies.

This work demonstrates that achieving stable RL fine-tuning requires balancing optimization aggressiveness with reward clarity and model capacity. Future research should extend these experiments to a broader range of model sizes within the same architecture families to provide deeper insights into scaling effects. Further, systematic analysis of hyperparameters alongside richer reward functions could significantly advance the stability, efficiency, and reliability of RL-based LLM fine-tuning for retrieval and structured generation tasks.